\subsection{Single Transferable Vote (STV)}

STV is a ranked-choice proportional system in which voters rank individual
candidates from all parties in preference order.  Seats are filled through an
iterative elimination-and-transfer process \citep{gallagher2011ireland}:

\begin{enumerate}
  \item Compute the Droop quota: $Q = \lfloor P / (S + 1) \rfloor + 1$,
        where $P$ is total valid ballots and $S$ is the number of seats.
  \item Candidates who reach quota are elected; their surplus votes
        (above quota) are redistributed to next preferences at a
        transfer value $TV = \text{surplus} / \text{total votes for
        candidate}$.
  \item If no candidate reaches quota, eliminate the candidate with
        the fewest votes and redistribute their votes at full value.
  \item Repeat until all seats are filled.
\end{enumerate}

STV produces proportional outcomes when parties are cohesive in their voter
rankings.  Its main advantage over closed-list PR is that voters can express
preferences within parties, not just between parties.  Its main disadvantage
is ballot complexity: in a 4-member district with 5 parties each running 2
candidates, a ballot may rank up to 10 candidates.

STV is used in Ireland (Dail Eireann, 3--5 member constituencies), Malta
(5 members), Scotland (local government, 3--4 members), Australia (Senate,
state-by-state, 6--12 members), and several U.S. and Canadian cities.  The
Irish system, with 3--5 member constituencies, is the closest international
model to a potential U.S. House reform.

\subsection{Open-List Proportional Representation (OLPR)}

OLPR is a party-list system in which voters choose a party and may
additionally cast a candidate preference vote \citep{karvonen2004preferential}.
Seats are allocated to parties using a divisor method (D'Hondt or
Sainte-Lague), and within each party, seats go to candidates with the most
preference votes.  If a voter does not express a candidate preference, their
vote counts toward the party's total but not toward any specific candidate.

OLPR is used in Sweden, Finland, the Netherlands, and Brazil (Chamber of
Deputies).  The defining feature is that voters can influence which
individuals from their preferred party are elected, while the party-list
structure ensures that party seat shares are proportional to vote shares.

For algorithmic implementation, OLPR is simpler than STV: the seat-allocation
step is a D'Hondt divisor calculation, well-established in the literature
\citep{balinski1982fair}.  The candidate-preference ranking within parties is
a separate computation that we implement as a post-processing step.

\subsection{District Magnitude and Compactness}

District magnitude---the number of seats per district---is the central design
variable in MMD systems.  Higher magnitude improves proportionality (lower
Gallagher index) but tends to produce less compact districts, because larger
population districts must span more territory.

The relationship between magnitude and compactness has been studied
theoretically \citep{taagepera2007seats} but rarely empirically for the
United States.  \citet{altman2011bqap} noted that algorithmic redistricting
produces more compact districts as $k$ increases (more, smaller districts),
which is the opposite direction from MMD design (fewer, larger districts).

For a fixed total population $P$ and $k$ districts of magnitude $M$ (so
$kM = S$ total seats), the mean district population is $P / k$.  A district
with twice the population of a single-member district is not necessarily twice
as spatially spread, because population is concentrated---but in low-density
states, larger districts do span more land, increasing perimeter and reducing
Polsby-Popper.

\subsection{MMDs and Minority Representation}

Under single-member districts, the primary mechanism for minority
representation is the majority-minority district (MMD): a district where a
minority group comprises $>50\%$ of the voting-age population, enabling
the election of minority-preferred candidates \citep{thernstrom1987whose}.
The Voting Rights Act (Section 2) can require such districts when specific
legal conditions are met \citep{thornburg1986}.

Under multi-member proportional districts, minority representation can
occur without geographic concentration.  A district where a minority group
is 35\% of the population electing 4 representatives would typically elect
1--2 minority-preferred candidates under proportional allocation.  This
``representation without packing'' was argued by \citet{guinier1994tyranny}
as a superior approach to the geographic concentration required by
majority-minority single-member districts.

The tradeoff is that MMD minority representation is contingent on minority
voting cohesion and the availability of minority candidates across party
lists.  Geographic concentration in a single-member district guarantees
minority representation; proportional allocation in an MMD converts minority
vote share into seat share only when the minority votes cohesively for
candidates of their choosing.
